% ---
% id: receiver/lem:fock-permanent-maps
% kind: lem
% title: "The lift of a map between two mode sets"
% summary: >-
%   The Fock lift of a linear map between two mode sets respects composition
%   and adjoints and keeps an isometry an isometry, and its entries in the
%   occupation bases are permanents divided by square roots of factorials.
%   Reading a block through a port dilation needs both.
% prominence: technical
% status: implemented
% uses: [def:fock, lem:fock-permanent]
% mentions: [lem:fock-functor, thm:factor-blocks]
% assumes: []
% refs: []
% provenance: newly-derived
% implements: certiqs_key.receiver.factor:_column
% lean:
% notes: "Kept apart from lem:fock-permanent, which is stated for an operator on one mode set, so that the claims resting on that lemma keep their review scope; the two-space case reads the same words and factorials in two mode sets."
% ---

Lemma~\ref{lem:fock-functor} and Lemma~\ref{lem:fock-permanent} are stated for
an operator on one mode space. Theorem~\ref{thm:factor-blocks} lifts an
isometry from the declared inputs into a different set of modes, and needs the
same two facts for it: that the lift composes and takes adjoints, and that its
entries are permanents.

\begin{lemma}\label{lem:fock-permanent-maps}
Let $I$, $O$ and $O'$ be finite mode sets, each with an orthonormal basis
listed in a fixed order, let $X:\mathbb{C}^{I}\to\mathbb{C}^{O}$ and
$X':\mathbb{C}^{O}\to\mathbb{C}^{O'}$ be linear maps, and let $\Fock(X)$ be
the lift of Definition~\ref{def:fock}.
\begin{enumerate}[label=(\roman*),leftmargin=3.0em]
\item $\Fock(X)$ maps $\Symn{n}(\mathbb{C}^{I})$ into
  $\Symn{n}(\mathbb{C}^{O})$ for every $n\ge0$;
  $\Fock(X'X)=\Fock(X')\Fock(X)$ and $\Fock(X)^{\dagger}=\Fock(X^{\dagger})$;
  and if $X^{\dagger}X=\id_{I}$, then $\Fock(X)^{\dagger}\Fock(X)$ is the
  identity of $\mathcal{F}(\mathbb{C}^{I})$, so that $\Fock(X)$ is an
  isometry.
\item Let $\vec m\in\mathbb{N}_0^{O}$ and $\vec m'\in\mathbb{N}_0^{I}$ be
  occupation vectors with $|\vec m|=|\vec m'|=n$, let $w_{\vec m}$,
  $w_{\vec m'}$, $\vec m!$ and $\vec m'!$ be the words and the factorial
  products of Lemma~\ref{lem:fock-permanent}, each formed in its own mode set,
  and let $X[\vec m|\vec m']$ be the $n\times n$ matrix with entries
  $X_{w_{\vec m}(r),\,w_{\vec m'}(s)}$. Then
  \begin{equation}\label{eq:fock-permanent-maps}
    \bra{\vec m}\Fock(X)^{(n)}\ket{\vec m'}
    =\frac{\operatorname{per}X[\vec m|\vec m']}{\sqrt{\vec m!\,\vec m'!}} ,
  \end{equation}
  with $\operatorname{per}$ as in Lemma~\ref{lem:fock-permanent}; at $n=0$
  both sides are $1$.
\end{enumerate}
\end{lemma}

\begin{proof}
(i) On $\Symn{n}(\mathbb{C}^{I})$ the lift $\Fock(X)$ is the restriction of
$X^{\otimes n}$, which carries $\Symn{n}(\mathbb{C}^{I})$ into
$\Symn{n}(\mathbb{C}^{O})$ (Definition~\ref{def:fock}); so $\Fock(X)$ maps
each photon-number summand into the summand of the same $n$. Factor by
factor, $(X'X)^{\otimes n}=X'^{\otimes n}X^{\otimes n}$, and since
$X^{\otimes n}$ lands in $\Symn{n}(\mathbb{C}^{O})$, the restriction of the
product to $\Symn{n}(\mathbb{C}^{I})$ is the product of the restrictions;
summing over $n$ gives $\Fock(X'X)=\Fock(X')\Fock(X)$. On the full tensor
powers $(X^{\otimes n})^{\dagger}=(X^{\dagger})^{\otimes n}$, and
$(X^{\dagger})^{\otimes n}$ carries $\Symn{n}(\mathbb{C}^{O})$ into
$\Symn{n}(\mathbb{C}^{I})$ (Definition~\ref{def:fock} for the map
$X^{\dagger}$). So for $\psi\in\Symn{n}(\mathbb{C}^{O})$ and
$\phi\in\Symn{n}(\mathbb{C}^{I})$,
$\langle\psi,X^{\otimes n}\phi\rangle=\langle(X^{\dagger})^{\otimes n}\psi,\phi\rangle$
with both vectors on the right in their symmetric subspaces, and the
restriction of $(X^{\dagger})^{\otimes n}$ to $\Symn{n}(\mathbb{C}^{O})$ is
the adjoint of the restriction of $X^{\otimes n}$ to
$\Symn{n}(\mathbb{C}^{I})$. Distinct $n$ give orthogonal summands, so
$\Fock(X)^{\dagger}=\Fock(X^{\dagger})$. If $X^{\dagger}X=\id_{I}$, these give
$\Fock(X)^{\dagger}\Fock(X)=\Fock(X^{\dagger}X)=\Fock(\id_{I})$, and
$\id^{\otimes n}$ restricts to the identity of every
$\Symn{n}(\mathbb{C}^{I})$.

(ii) Let $P_I=\frac1{n!}\sum_{\pi}\Pi^{I}_{\pi}$ and
$P_O=\frac1{n!}\sum_{\pi}\Pi^{O}_{\pi}$ be the orthogonal projectors onto
$\Symn{n}(\mathbb{C}^{I})$ and $\Symn{n}(\mathbb{C}^{O})$, $\Pi^{I}_{\pi}$ and
$\Pi^{O}_{\pi}$ permuting the tensor factors, and let $\ket{w_{\vec m}}$ and
$\ket{w_{\vec m'}}$ be the product basis vectors whose $r$-th factors are the
basis vectors of the letters $w_{\vec m}(r)$ and $w_{\vec m'}(r)$. The
permutations that fix $\ket{w_{\vec m}}$ are those that permute equal letters
among themselves, $\vec m!$ of them, and every other $\Pi^{O}_{\pi}$ sends it
to an orthogonal product basis vector; so
$\lVert P_O\ket{w_{\vec m}}\rVert^{2}=\bra{w_{\vec m}}P_O\ket{w_{\vec m}}
=\vec m!/n!$, and the symmetrized and normalized vector of
Definition~\ref{def:fock} is
$\ket{\vec m}=\sqrt{n!/\vec m!}\,P_O\ket{w_{\vec m}}$, and likewise
$\ket{\vec m'}=\sqrt{n!/\vec m'!}\,P_I\ket{w_{\vec m'}}$. For every
permutation $\pi$, $X^{\otimes n}\Pi^{I}_{\pi}=\Pi^{O}_{\pi}X^{\otimes n}$,
both sending $v_1\otimes\cdots\otimes v_n$ to the tensor product of the
$Xv_r$ permuted by $\pi$; so $X^{\otimes n}P_I=P_OX^{\otimes n}$. With
$P_O^{\dagger}=P_O$ and $P_I^{2}=P_I$,
\begin{equation*}
  \bra{\vec m}X^{\otimes n}\ket{\vec m'}
  =\frac{n!}{\sqrt{\vec m!\,\vec m'!}}
   \bra{w_{\vec m}}P_OX^{\otimes n}P_I\ket{w_{\vec m'}}
  =\frac{n!}{\sqrt{\vec m!\,\vec m'!}}
   \bra{w_{\vec m}}X^{\otimes n}P_I\ket{w_{\vec m'}}
  =\frac{1}{\sqrt{\vec m!\,\vec m'!}}\sum_{\pi}
   \bra{w_{\vec m}}X^{\otimes n}\Pi^{I}_{\pi}\ket{w_{\vec m'}} .
\end{equation*}
The vector $\Pi^{I}_{\pi}\ket{w_{\vec m'}}$ is the product vector whose $r$-th
factor is the basis vector of $w_{\vec m'}(\sigma(r))$ for a permutation
$\sigma$ determined by $\pi$, and $\pi\mapsto\sigma$ is a bijection of the
permutation group. The summand is therefore
$\prod_{r}X_{w_{\vec m}(r),\,w_{\vec m'}(\sigma(r))}$, and the sum over
$\sigma$ is $\operatorname{per}X[\vec m|\vec m']$. The left-hand side is the
entry of $\Fock(X)^{(n)}$, $\ket{\vec m'}$ lying in
$\Symn{n}(\mathbb{C}^{I})$. At $n=0$ the lift is $1$ on the vacuum, the
permanent of the empty matrix is $1$ and both factorial products are $1$.
\end{proof}
